%%%%%%%%%%%%%%%%%%%%%%%%%%%%%%%%%%%%%%%%%%%%%%%%%%%%%%%%%%%%%%%%%%%%%%%%%%%%%%
% Much of this code was taken from Charlie's maths template, available here: %
% https://github.com/SeniorMars/dotfiles/tree/main/latex_template            %
%%%%%%%%%%%%%%%%%%%%%%%%%%%%%%%%%%%%%%%%%%%%%%%%%%%%%%%%%%%%%%%%%%%%%%%%%%%%%%


%================================
% COLORS
%================================

\definecolor{doc}{RGB}{0,60,110}
\definecolor{myg}{RGB}{56, 140, 70}        % claim accent
\definecolor{myp}{RGB}{197, 92, 212}       % corollary accent
\definecolor{myb}{RGB}{45, 111, 177}       % question accent

\definecolor{mytheorembg}{HTML}{F2F2F9}
\definecolor{mytheoremfr}{HTML}{00007B}
\definecolor{mylemmabg}{HTML}{FFFAF8}
\definecolor{mylemmafr}{HTML}{983b0f}
\definecolor{myexamplebg}{HTML}{F2FBF8}
\definecolor{myexamplefr}{HTML}{88D6D1}
\definecolor{myexampleti}{HTML}{2A7F7F}
\definecolor{mydefinitbg}{HTML}{FFF5F5}
\definecolor{mydefinitfr}{HTML}{8A1C1C}
\definecolor{myremarkbg}{HTML}{FAF9F9}
\definecolor{myremarkfr}{HTML}{000000}

\definecolor{codebg}{HTML}{F7F8FA}
\definecolor{codegutter}{HTML}{EDEFF3}
\definecolor{codeframe}{HTML}{D8DEE8}
\definecolor{codeaccent}{HTML}{3867D6}
\definecolor{codekeyword}{HTML}{7C3AED}
\definecolor{codestring}{HTML}{0F766E}
\definecolor{codecomment}{HTML}{64748B}
\definecolor{codenumber}{HTML}{94A3B8}
\definecolor{codeidentifier}{HTML}{1F2937}

%================================
% TCOLORBOX BASE STYLES
%================================

\tcbuselibrary{theorems, skins, hooks, listings}

\tcbset{
	cmtleftpanel/.style 2 args={
		enhanced, breakable,
		colback = #1,
		frame hidden,
		boxrule = 0sp,
		borderline west = {2pt}{0pt}{#2},
		sharp corners,
		description font = \mdseries,
		separator sign none,
	},
	cmtleftthm/.style 2 args={
		cmtleftpanel={#1}{#2},
		detach title,
		before upper = \tcbtitle\par\smallskip,
		coltitle = #2,
		fonttitle = \bfseries\sffamily,
		segmentation style={solid, #2},
	},
	cmtexamplethm/.style={
		colback = myexamplebg,
		breakable,
		colframe = myexamplefr,
		coltitle = myexampleti,
		boxrule = 1pt,
		sharp corners,
		detach title,
		before upper=\tcbtitle\par\smallskip,
		fonttitle = \bfseries,
		description font = \mdseries,
		separator sign none,
		description delimiters parenthesis,
	},
	cmtdefinitionthm/.style={
		enhanced,
		before skip=2mm,
		after skip=2mm,
		colback=mydefinitbg,
		colframe=mydefinitfr,
		boxrule=0.5mm,
		attach boxed title to top left={xshift=1cm,yshift*=1mm-\tcboxedtitleheight},
		varwidth boxed title*=-3cm,
		colbacktitle=mydefinitfr,
		coltitle=white,
		boxed title style={
			frame code={
				\path[fill=tcbcolback]
				([yshift=-1mm,xshift=-1mm]frame.north west)
				arc[start angle=0,end angle=180,radius=1mm]
				([yshift=-1mm,xshift=1mm]frame.north east)
				arc[start angle=180,end angle=0,radius=1mm];
				\path[left color=tcbcolback!60!black,right color=tcbcolback!60!black,
					middle color=tcbcolback!80!black]
				([xshift=-2mm]frame.north west) -- ([xshift=2mm]frame.north east)
				[rounded corners=1mm]-- ([xshift=1mm,yshift=-1mm]frame.north east)
				-- (frame.south east) -- (frame.south west)
				-- ([xshift=-1mm,yshift=-1mm]frame.north west)
				[sharp corners]-- cycle;
			},
			interior engine=empty,
		},
		fonttitle=\bfseries,
	},
}

\makeatletter
\tcbset{
	cmtribbonbox/.style={
		enhanced,
		breakable,
		colback=white,
		attach boxed title to top left={yshift*=-\tcboxedtitleheight},
		fonttitle=\bfseries,
		boxed title size=title,
		boxed title style={
			sharp corners,
			rounded corners=northwest,
			colback=tcbcolframe,
			boxrule=0pt,
		},
		underlay boxed title={
			\path[fill=tcbcolframe] (title.south west)--(title.south east)
			to[out=0, in=180] ([xshift=5mm]title.east)--
			(title.center-|frame.east)
			[rounded corners=\kvtcb@arc] |-
			(frame.north) -| cycle;
		},
	},
}
\makeatother

%================================
% THEOREM-FAMILY BOXES
%================================

\newtcbtheorem[number within=section,
	crefname={theorem}{theorems}, Crefname={Theorem}{Theorems}]
	{theorem}{Theorem}{cmtleftthm={mytheorembg}{mytheoremfr}}{th}

\newtcbtheorem[number within=section,
	crefname={corollary}{corollaries}, Crefname={Corollary}{Corollaries}]
	{corollary}{Corollary}{cmtleftthm={myp!10}{myp!85!black}}{th}

\newtcbtheorem[number within=section,
	crefname={lemma}{lemmas}, Crefname={Lemma}{Lemmas}]
	{lemma}{Lemma}{cmtleftthm={mylemmabg}{mylemmafr}}{th}

\newtcbtheorem[number within=section,
	crefname={claim}{claims}, Crefname={Claim}{Claims}]
	{claim}{Claim}{cmtleftthm={myg!10}{myg!85!black}, borderline west={2pt}{0pt}{myg}}{th}

\newtcbtheorem[number within=section,
	crefname={example}{examples}, Crefname={Example}{Examples}]
	{example}{Example}{cmtexamplethm}{ex}

\newtcbtheorem[number within=section,
	crefname={definition}{definitions}, Crefname={Definition}{Definitions}]
	{definition}{Definition}{cmtdefinitionthm, title={#2}, #1}{def}

\newtcbtheorem[number within=section,
	crefname={remark}{remarks}, Crefname={Remark}{Remarks}]
	{remark}{Remark}{cmtleftthm={myremarkbg}{myremarkfr}}{rmk}

\newtcbtheorem[number within=section,
	crefname={question}{questions}, Crefname={Question}{Questions}]
	{question}{Question}{cmtribbonbox, colframe=myb!80!black, title={#2}, #1}{qst}

%================================
% NOTE BOX (unnumbered)
%================================

\usetikzlibrary{arrows, calc, shadows.blur}

\newtcolorbox{quotebox}[1][]{
	cmtleftpanel={myp!10}{myp!85!black},
	fontupper=\itshape,
	#1
}
\newtcolorbox{note}[1][]{
	enhanced jigsaw,
	colback=gray!20!white,
	colframe=gray!80!black,
	size=small,
	boxrule=1pt,
	title=\textbf{Note:-},
	halign title=flush center,
	coltitle=black,
	breakable,
	drop shadow=black!50!white,
	attach boxed title to top left={xshift=1cm,yshift=-\tcboxedtitleheight/2,yshifttext=-\tcboxedtitleheight/2},
	minipage boxed title=1.5cm,
	boxed title style={
		colback=white,
		size=fbox,
		boxrule=1pt,
		boxsep=2pt,
		underlay={
			\coordinate (dotA) at ($(interior.west) + (-0.5pt,0)$);
			\coordinate (dotB) at ($(interior.east) + (0.5pt,0)$);
			\begin{scope}
				\clip (interior.north west) rectangle ([xshift=3ex]interior.east);
				\filldraw [white, blur shadow={shadow opacity=60, shadow yshift=-.75ex}, rounded corners=2pt]
					(interior.north west) rectangle (interior.south east);
			\end{scope}
			\begin{scope}[gray!80!black]
				\fill (dotA) circle (2pt);
				\fill (dotB) circle (2pt);
			\end{scope}
		},
	},
}

%================================
% CODEBLOCK
%================================

\lstdefinestyle{cmtcodebase}{
	basicstyle=\ttfamily\footnotesize,
	keywordstyle=\bfseries\color{codekeyword},
	commentstyle=\itshape\color{codecomment},
	stringstyle=\color{codestring},
	identifierstyle=\color{codeidentifier},
	numberstyle=\scriptsize\color{codenumber},
	showstringspaces=false,
	tabsize=4,
	columns=fullflexible,
	keepspaces=true,
	upquote=true,
	breaklines=true,
	breakatwhitespace=false,
	postbreak=\mbox{\textcolor{codenumber}{$\hookrightarrow$}\space},
}
\lstdefinestyle{cmtcode}{
	style=cmtcodebase,
	numbers=left,
	numbersep=10pt,
	xleftmargin=2.8em,
}

\tcbset{
	cmtcodebox/.style={
		enhanced, breakable, sharp corners,
		boxrule=0pt, frame hidden,
		colback=codebg, colframe=codeframe,
		coltitle=codeidentifier,
		fonttitle=\bfseries\sffamily\small,
		colbacktitle=codegutter,
		titlerule=0pt,
		toptitle=4pt, bottomtitle=4pt,
		lefttitle=8pt, righttitle=8pt,
		left=8pt, right=8pt, top=6pt, bottom=6pt,
		boxsep=0pt,
		before skip=\baselineskip, after skip=\baselineskip,
		borderline west={2pt}{0pt}{codeaccent},
		segmentation hidden,
		listing only,
	},
}
\newtcblisting{codeblock}[2][]{
	cmtcodebox,
	title={#2},
	listing options={style=cmtcode, language=#2},
	#1
}

%================================
% SHORT COMMANDS
%================================

\newcommand{\dfn}[2]{\begin{definition}{#1}{}#2\end{definition}}
\newcommand{\ex}[2]{\begin{example}{#1}{}#2\end{example}}
\newcommand{\rmk}[2]{\begin{remark}{#1}{}#2\end{remark}}
\newcommand{\clm}[3]{\begin{claim}{#1}{#2}#3\end{claim}}
\newcommand{\thm}[2]{\begin{theorem}{#1}{}#2\end{theorem}}
\newcommand{\cor}[2]{\begin{corollary}{#1}{}#2\end{corollary}}
\newcommand{\mlemma}[2]{\begin{lemma}{#1}{}#2\end{lemma}}
\newcommand{\qst}[2]{\begin{question}{#1}{}#2\end{question}}
\newcommand{\nt}[1]{\begin{note}#1\end{note}}
\newcommand{\qte}[1]{\begin{quotebox}#1\end{quotebox}}
